\begingroup
\setlength{\LTcapwidth}{\textwidth}
\begin{longtable}{@{}p{1.2cm}p{8.4cm}p{1.55cm}p{1.45cm}@{}}
\caption{Kebutuhan fungsional sistem.}
\label{tab:tab3.1_kebutuhan_fungsional} \\
\toprule
\textbf{Kode} & \textbf{Kebutuhan Fungsional} & \textbf{Asal} & \textbf{Tujuan} \\
\midrule
\endfirsthead

\multicolumn{4}{@{}l}{\small\emph{Tabel \thetable{} (lanjutan)}} \\
\toprule
\textbf{Kode} & \textbf{Kebutuhan Fungsional} & \textbf{Asal} & \textbf{Tujuan} \\
\midrule
\endhead

\bottomrule
\endlastfoot

KF-01 & Sistem menerima masukan catatan \emph{logbook} terstruktur yang mencakup kadar glukosa, asupan karbohidrat, pemberian insulin, dan aktivitas fisik. & M2 & T1 \\[2pt]

KF-02 & Sistem menghitung besaran turunan berbasis fisiologi dari catatan mentah, mencakup \emph{insulin on board}, \emph{carbohydrate on board}, laju perubahan glukosa, dan penyandian siklis waktu. & M2 & T1 \\[2pt]

KF-03 & Sistem memprediksi kadar glukosa pada horizon 30 menit dan 60 menit beserta selang ketidakpastiannya. & M1, M2 & T1 \\[2pt]

KF-04 & Sistem menyusun kueri penelusuran yang dikondisikan pada kadar glukosa terprediksi, bukan pada kadar glukosa yang sedang berlaku. & M3 & T2 \\[2pt]

KF-05 & Sistem menelusuri korpus pedoman klinis dan mengembalikan potongan dokumen yang relevan beserta penunjuk dokumen dan nomor halaman sumbernya. & M3 & T2 \\[2pt]

KF-06 & Sistem menyusun rekomendasi klinis berbahasa Indonesia yang dibumikan pada potongan dokumen hasil penelusuran. & M2, M3 & T2 \\[2pt]

KF-07 & Sistem menyajikan kondisi klinis terstruktur berupa tingkat risiko, arah perubahan, dan tingkat kegentingan, serta memberikan peringatan dini hipoglikemia dan
peringatan divergensi. Peringatan divergensi diberikan ketika kondisi terprediksi berbeda kategori dari kondisi yang sedang berlaku, misalnya kadar glukosa terkini berada
dalam rentang sasaran sedangkan kondisi terprediksi tergolong hipoglikemia. & M1, M2 & T1, T2 \\[2pt]

KF-08 & Sistem mendukung alur \emph{doctor-mediated}, yaitu sistem tidak melakukan tindakan otonom dan rekomendasi disajikan sebagai bahan pertimbangan yang dapat diperiksa,
sehingga keputusan klinis akhir tetap berada pada dokter. & M4 & T3 \\
\end{longtable}

\noindent\tabnotestyle Kolom asal mengacu pada persoalan yang diuraikan pada Subbab~\ref{subsec:bab3-identifikasi-persoalan}, yaitu M1 pengelolaan bersifat reaktif, M2 penafsiran manual dan tidak antisipatif, M3 prediksi terputus dari rekomendasi, dan M4 solusi berinfrastruktur berat tidak dapat diterapkan. Kolom tujuan mengacu pada tujuan penelitian pada Subbab~\ref{sec:bab1-tujuan}, yaitu T1 modul prediksi kadar glukosa, T2 mekanisme RAG dengan \emph{query transformation} terkondisi-prediksi, dan T3 integrasi keduanya ke dalam CDSS beralur \emph{doctor-mediated}.\normalsize
\endgroup
